% TOP_PROBLEM: {"id": 30001375, "title": "Permanent–Determinant Separation Through Geometric Complexity Theory", "category_id": 15, "category": {"id": 15, "name": "computer_science", "display_name": "Computer Science", "description": "Computational complexity, algorithms, and theoretical CS.", "slug": "computer-science", "order_index": 15, "created_at": "2026-07-31T15:26:25.671Z"}, "status": "open"}
% FIRST_POSTED: 2026-09-27
% ATTEMPT_STATUS: unresolved
% !TeX program = lualatex
\documentclass[11pt]{article}
\usepackage[margin=25mm]{geometry}
\usepackage{fontspec}
\setmainfont{Latin Modern Roman}
\usepackage{amsmath,amssymb,mathtools}
\usepackage{unicode-math}
\setmathfont{Latin Modern Math}
\usepackage{xurl,hyperref}
\hypersetup{colorlinks=true,urlcolor=blue}
\setcounter{secnumdepth}{0}
\setlength{\parindent}{0pt}
\setlength{\parskip}{5pt}
\setlength{\emergencystretch}{3em}
\begin{document}
\section*{\texorpdfstring{Permanent–Determinant Separation Through Geometric Complexity Theory}{Permanent–Determinant Separation Through Geometric Complexity Theory}}\label{30001375}
\subsection{Definitions and mathematical statement}
For variables $x_{ij}$ over a characteristic-zero field $K$, define
\[
 \operatorname{per}_n(x)=\sum_{\sigma\in S_n}\prod_{i=1}^n x_{i,\sigma(i)}.
\]
An algebraic circuit computes a polynomial using field constants, variables, addition, and multiplication. The class $\mathsf{VP}$ has polynomially bounded numbers of variables, degrees, and circuit sizes; $\mathsf{VNP}$ consists of polynomially long Boolean sums of such efficiently computable polynomials. The selected claim is that no polynomial bounds the arithmetic circuit sizes of $\operatorname{per}_n$, equivalently $\mathsf{VP}\ne\mathsf{VNP}$ in characteristic zero. This is a nonuniform algebraic-operation model, not a bit-complexity assertion.
\par\smallskip\textit{Formulation and concept references:} \href{https://www.prateekdwivedi.in/papers/symACT-confversion.pdf}{[S1]}.

\subsection{Short English statement}
Can one prove that the permanent requires more than polynomially many arithmetic operations in general algebraic circuits?
\subsection{Research attempt}

\paragraph{Scope, sources, and the chosen attack (27 September 2026).}
Fix the characteristic-zero field $K$. Constants from $K$ are allowed,
and the circuit for one matrix size need not be computably recoverable
from the circuit for another. The desired conclusion is that no
polynomial bounds the minimum sizes for all $n$; a lower bound against
one restricted circuit form does not suffice.
The symmetric classes in [S1, S2] impose conditions absent here, and
[S1] explicitly distinguishes their orbit-size measure from general
circuit size. The August 2026 report [S3] gives elusive-function and
bounded-depth results; its discussion identifies degree growth as an
obstruction to deriving the full separation. None is read as a solution
to unrestricted $\mathsf{VP}\ne\mathsf{VNP}$.

The strategy is to calculate concrete derivative invariants of the
permanent, extract the lower bounds they actually prove, and test their
ability to survive general multiplication gates. A second calculation
recovers the known quadratic determinantal bound of Mignon--Ressayre
[E2]. This is a reconstruction and stress test of classical tools, not
a claim of a new circuit lower bound. Valiant's completeness theorem
[E1] supplies the stated equivalence with the class separation; the
membership direction can be proved explicitly here.

\paragraph{An explicit Boolean-sum witness and an exact upper bound.}
For $y=(y_1,\ldots,y_n)$ define
\begin{equation}\label{vp:witness}
 g_n(x,y)=(-1)^n\prod_{j=1}^n(1-2y_j)
                   \prod_{i=1}^n\left(\sum_{j=1}^n x_{ij}y_j\right).
\end{equation}
This polynomial has degree at most $3n$, $n^2+n$ variables, and a
circuit of size $O(n^2)$. For Boolean $y$ with support $S$,
its value is $(-1)^{n-|S|}\prod_i\sum_{j\in S}x_{ij}$. Expanding the
sum of these values over all $S$ gives, for each function
$f:[n]\to[n]$, the monomial $\prod_i x_{i,f(i)}$ with coefficient
\[
 \sum_{S\supseteq\operatorname{im}f}(-1)^{n-|S|}
 =\begin{cases}1,&\operatorname{im}f=[n],\\0,&\text{otherwise}.
 \end{cases}
\]
The identity follows by choosing the subset of the complement; a
nonempty complement contributes $(1-1)^d=0$. A surjective map between
these two $n$-element sets is a permutation. Therefore
\begin{equation}\label{vp:boolean}
 \operatorname{per}_n(x)=\sum_{y\in\{0,1\}^n}g_n(x,y).
\end{equation}
This proves the required $\mathsf{VNP}$ membership with polynomially
many witness bits. It also gives an exact $O(n^2 2^n)$ circuit by
explicit summation, the usual inclusion--exclusion mechanism.
The phrase ``polynomially long Boolean sum'' bounds the number of
witness bits, not the number of summands. Evaluating $g_n$ efficiently
does not make the exponential summation a polynomial-size computation.

For $n=1$, the two Boolean evaluations give $x_{11}$. For $n=2$,
the subset terms cancel the two monomials that reuse a column and leave
$x_{11}x_{22}+x_{12}x_{21}$. Thus the witness has the correct signs
even at the first nontrivial size. Direct enumeration for sizes up to
five checked the identity on integer matrices in the local calculation.

\paragraph{Exact derivative dimensions.}
For a homogeneous degree-$n$ polynomial $f$, let $D_k(f)$ be the
$K$-linear span of all order-$k$ partial derivatives, and set
$d_k(f)=\dim_K D_k(f)$. Here $0\le k\le n$.
A nonzero derivative of $\operatorname{per}_n$ differentiates distinct
variables occupying distinct rows and columns. It is exactly the
permanent on the undeleted rows and columns. Conversely every choice
of $k$ deleted rows and $k$ deleted columns occurs by choosing a
matching between them. Hence
\begin{equation}\label{vp:derivatives}
 d_k(\operatorname{per}_n)=\binom nk^2.
\end{equation}
To justify equality rather than just a spanning bound, observe that
the surviving monomials determine their row set and column set.
Different subpermanents consequently have disjoint monomial supports.
They are nonzero and linearly independent. At $k=n$ there is just the
constant polynomial, agreeing with the formula.

This yields an exponential lower bound for a genuine restricted model.
Suppose
\[
 \operatorname{per}_n=\sum_{a=1}^s c_a\ell_a^n
\]
with homogeneous linear forms $\ell_a$ and scalars $c_a$. Every
order-$k$ derivative belongs to the span of the $s$ polynomials
$\ell_a^{n-k}$. Thus
\[
 s\ge\binom n{\lfloor n/2\rfloor}^{\!2}
       \ge\frac{4^n}{(n+1)^2}.
\]
The last inequality uses the largest of the $n+1$ binomial
coefficients and their sum $2^n$. This is a lower bound for sums of
powers of linear forms, not for unrestricted depth-three circuits or
general circuits. No reduction preserving polynomial size from the
latter models to this restricted representation has been proved here.

\paragraph{Why derivative dimension does not control general circuits.}
The elementary product $h(z)=z_1\cdots z_n$ has a circuit with $n-1$
multiplications, but $d_k(h)=\binom nk$, because its derivatives are
all squarefree monomials of degree $n-k$. This already refutes the
proposal that exponential derivative dimension implies exponential
general circuit size. The product rule explains the failure:
\[
 D_k(fh)\subseteq\sum_{j=0}^k D_j(f)D_{k-j}(h),\qquad
 d_k(fh)\le\sum_{j=0}^k d_j(f)d_{k-j}(h),
\]
where a product of spaces denotes the span of pairwise products.
Repeated multiplication can generate binomially many derivatives
without using binomially many gates.

There is a stronger countertest: the determinant has exactly the same
entire derivative-dimension profile as the permanent. Its nonzero
order-$k$ derivatives are signed complementary minors, which are
independent by the same disjoint-support argument. Thus
\[
 d_k(\det_n)=\binom nk^2=d_k(\operatorname{per}_n)
 \quad(0\le k\le n).
\]
Nevertheless $\det_n$ has polynomial-size circuits over $K$. For
completeness, compute $p_j=\operatorname{tr}(X^j)$ for $j\le n$ by
successive matrix products, at cost $O(n^4)$, and use
\[
 e_0=1,\qquad
 e_k=\frac1k\sum_{j=1}^k(-1)^{j-1}e_{k-j}p_j.
\]
Newton's identities give $e_n=\det X$. These identities follow by
logarithmically differentiating
$\prod_i(1+\lambda_i t)$ and comparing coefficients; polynomial
identity then gives them for the generic matrix. The divisions by
the positive integers $k$ are permitted field constants in
characteristic zero, not divisions by input expressions.
Consequently no invariant that depends only on this derivative-rank
profile can distinguish the two families. A successful refinement
must retain additional structure or exploit restrictions on gates.

\paragraph{A more local calculation: the quadratic determinant bound.}
Let $\operatorname{dc}(f)$ be the least $m$ for which
$f(x)=\det_m(A_0+\sum x_{ij}A_{ij})$ over $K$.
If $H_f$ denotes the Hessian, the affine chain rule gives
\begin{equation}\label{vp:hessian}
 H_f(x_0)=L^{\mathsf T}H_{\det_m}(A(x_0))L.
\end{equation}
At a singular $m$ by $m$ matrix the determinant Hessian has rank at
most $2m$. Row and column transformations reduce the matrix to
$\operatorname{diag}(I_r,0)$. For $r\le m-3$ its quadratic term
vanishes. For $r=m-2$ it is a two-by-two determinant, of Hessian
rank four. For $r=m-1$ the quadratic term is
\[
 z_{mm}\sum_{i<m}z_{ii}-\sum_{i<m}z_{im}z_{mi},
\]
whose Hessian rank is at most $2+2(m-1)=2m$.
The case $m=1$ has zero Hessian.

Here is an explicit permanent calculation to apply that bound.
Take $X_0$ with all entries one except $(X_0)_{11}=1-n$.
Then $\operatorname{per}_n(X_0)=(n-1)!(n-1+1-n)=0$.
For $n\ge3$, divide its Hessian by $(n-3)!$.
An entry for two variables in a common row or column is zero;
for compatible variables it is $t=n-2$ if their deleted rows or
columns include index one, and $-2$ otherwise.

To check nonsingularity directly, write a kernel vector in blocks
$(a,(b_j),(c_i),(d_{ij}))$ for the corner, first row, first column,
and interior. All interior indices range over $2,\ldots,n$.
Write $S_b,S_c,S_d$ for the corresponding sums, and $R_i,C_j$ for
interior row and column sums. The corner and border equations give
\[
 tS_d=0,\quad t(S_c+S_d-C_j)=0,\quad
 t(S_b+S_d-R_i)=0.
\]
Summing them yields $S_b=S_c=S_d=R_i=C_j=0$.
The interior equations reduce to
\[
 2d_{ij}=t(a-b_j-c_i).
\]
Their row sums give $c_i=a$, their column sums give $b_j=a$, and
$S_b=0$ forces $a=0$. All coordinates vanish. For $n=2$, the
Hessian of $x_{11}x_{22}+x_{12}x_{21}$ has rank four directly.
Thus $\operatorname{rank}H_{\operatorname{per}_n}(X_0)=n^2$ for
$n\ge2$. Substitution in \eqref{vp:hessian} proves the classical
Mignon--Ressayre bound [E2]
\begin{equation}\label{vp:dc}
 \operatorname{dc}(\operatorname{per}_n)\ge n^2/2.
\end{equation}

\paragraph{What that bound leaves open.}
The Hessian argument is stronger than merely comparing the global
dimensions $d_k$: it evaluates derivatives at a specially chosen zero
and uses the geometry of singular determinant matrices. It still
produces a polynomial lower bound, compatible with determinantal
representations of size $n^3$, for example. Its input Hessian has only
$n^2$ rows, so comparing its rank with $2m$ cannot by itself produce
a superquadratic lower bound in $n$.

There is a second, independent gap in using this route for the selected
problem. A polynomial-size affine determinantal representation would
give a polynomial-size general circuit, by the determinant algorithm
above. The reverse implication is not known with polynomial overhead:
determinant is complete for polynomial-size algebraic branching
programs, a potentially smaller model than $\mathsf{VP}$ [E3,
Sections 2.5 and 2.9]. Thus even a superpolynomial lower bound on
$\operatorname{dc}(\operatorname{per}_n)$ would not automatically
separate $\mathsf{VP}$ from $\mathsf{VNP}$. It would do so under the
additional hypothesis that every polynomial-size, polynomial-degree
general circuit family admits polynomial-size affine determinantal
representations. That simulation is an extra hypothesis, not a step
established by \eqref{vp:hessian}.

\paragraph{Why a generic counting argument does not select the permanent.}
There are indeed degree-$n$ polynomials in $n^2$ variables with large
circuits. To see the distinction from this task, fix a size bound $s$
for binary arithmetic circuits. There are finitely many circuit
topologies and choices of variable labels of that size. A topology
has at most $2s$ constant input slots; regard them as parameters.
The coefficients of its homogeneous degree-$n$ component are
polynomial functions of those parameters. Their image in the
coefficient space of dimension
\[
 N=\binom{n^2+n-1}{n}
\]
has Zariski closure of dimension at most $2s$: the field generated
by these coefficient functions has transcendence degree at most the
number of parameters. If $2s<N$, each such closure is proper.
A finite union of these proper closed sets cannot fill $K^N$ over
the infinite field $K$. Equivalently, the product of nonzero defining
polynomials for the closures is nonzero and cannot vanish at every
point of $K^N$. This proves existence of a polynomial beyond the
specified circuit bound, even while allowing arbitrary constants.

However, the permanent is one fixed coefficient vector, not a generic
one. The argument supplies no reason it avoids any of these circuit
images, and the chosen hard polynomial is not shown to have a
polynomial-length Boolean-sum description. Counting unrestricted
coefficient vectors therefore does not give a hard family in
$\mathsf{VNP}$. This also explains why counting circuit strings with
only a finite alphabet of constants would address a different model:
the actual constants range over all of $K$. The parameter argument
handles that issue but leaves the essential explicitness gap intact.

\paragraph{Verification, first gap, and outcome.}
Exact finite checks compared the Boolean witness sum with the
permutation definition, derivative-support counts with
$\binom nk^2$, and the displayed Hessian with direct complementary
subpermanent calculations. Hessian ranks were computed over the
rationals for $2\le n\le7$, yielding $n^2$ each time. The proof above,
not those finite ranks, handles arbitrary $n$. Characteristic zero
was used in Newton's identities, the kernel calculation, and the
completeness formulation; these arguments are not silently transferred
to characteristic two, where permanent and determinant coincide.

The first gap in the derivative strategy is a nontrivial upper bound
for a stronger permanent invariant on \emph{all} small circuits,
consistent with the product and determinant countertests. Three useful
next tasks are to preserve row/column structure in a refined derivative
space; to prove gate-by-gate bounds that account for reuse and
cancellation; and to identify a higher-order local obstruction whose
growth exceeds every fixed polynomial after the relevant simulation
cost is included. Merely enlarging the ordinary derivative spaces or
finding more small Hessian examples cannot meet these requirements.

\textbf{UNRESOLVED.} The attempt proves explicit membership and upper
bounds, an exponential lower bound for sums of powers, and reconstructs
the known quadratic determinantal lower bound. It does not give a
superpolynomial lower bound for general arithmetic circuits.
Written by GPT-6 Astra Ultra with primary-source scope checks and exact
finite calculations; no formal verification or novelty is asserted.

\subsection{Sources}
\begin{itemize}
\item[S1]Dwivedi, Pago, and Seppelt, Lower Bounds in Algebraic Complexity via Symmetry and Homomorphism Polynomials, STOC 2026, Introduction.
\url{https://www.prateekdwivedi.in/papers/symACT-confversion.pdf}.
\textit{Formulation source.}
\item[S2]Symmetric Algebraic Circuits and Homomorphism Polynomials.
\url{https://drops.dagstuhl.de/storage/00lipics/lipics-vol362-itcs2026/LIPIcs.ITCS.2026.46/LIPIcs.ITCS.2026.46.pdf}.
\textit{Further reference; not a proof endorsement.}
\item[S3]Arithmetic circuit lower bounds from sumset expansion.
\url{https://eccc.weizmann.ac.il/report/2026/138/}.
\textit{Further reference; not a proof endorsement.}
\item[E1]Leslie G. Valiant, \emph{Completeness classes in algebra}, Proceedings of STOC 1979, 249--261.
\url{https://doi.org/10.1145/800135.804419}.
Author-paper scan: \url{https://sites.math.washington.edu/~billey/colombia/references/valiant.completeness.1979.pdf}.
\textit{Primary source for algebraic completeness; the explicit Boolean-sum membership identity is proved above. Consulted 27 September 2026.}
\item[E2]Thierry Mignon and Nicolas Ressayre, \emph{A quadratic bound for the determinant and permanent problem}, International Mathematics Research Notices 2004(79), 4241--4253.
\url{https://math.univ-lyon1.fr/~ressayre/PDFs/permdet.pdf}.
\textit{Primary source of the known characteristic-zero quadratic determinantal lower bound, reconstructed here by an explicit Hessian calculation; consulted 27 September 2026.}
\item[E3]Peter B\"urgisser, \emph{Completeness classes in algebraic complexity theory}, arXiv:2406.06217 (2024), Sections 2.4--2.9.
\url{https://arxiv.org/html/2406.06217v1}.
\textit{Author account of the determinant, branching-program completeness, and the distinction from general circuits; consulted 27 September 2026.}
\end{itemize}
\end{document}
